\subsection{API interna y Route Handlers}
\label{subsec:internalApiRoutes}

En \texttt{Next.js}, los \textit{Route Handlers} permiten crear manejadores HTTP personalizados dentro del directorio \texttt{app}, usando las APIs web estandarizadas \texttt{Request} y \texttt{Response} \citep{nextjsroutehandlers2026}. En este proyecto esa capacidad se materializa en la carpeta \path{app/api/}, donde cada fichero \path{route.ts} expone operaciones de lectura o escritura mediante métodos HTTP como \texttt{GET}, \texttt{POST}, \texttt{PUT}, \texttt{PATCH} y \texttt{DELETE}, métodos soportados por la convención oficial de rutas de \texttt{Next.js} \citep{nextjsroutefile2026}. Desde el punto de vista HTTP, estos métodos expresan la intencion de la petición, por ejemplo recuperar un recurso, enviar datos, sustituirlos, modificarlos parcialmente o eliminarlos \citep{mdnhttpmethods2026}.

El uso de rutas API internas aporta varias ventajas al proyecto. En primer lugar, permite mantener en el mismo repositorio la interfaz y la frontera de servidor, reduciendo la distancia entre componentes visuales, contratos TypeScript y servicios de dominio. En segundo lugar, encaja con el patron \textit{Backend for Frontend}, recomendado por la propia documentación de \texttt{Next.js} para adaptar, filtrar o agregar datos antes de entregarlos a la interfaz \citep{nextjsbackendforfrontend2026}. En tercer lugar, evita exponer directamente la base de datos o los servicios internos al navegador: las páginas llaman a \path{/api/...}, mientras que los \textit{Route Handlers} validan sesión, rol, parámetros y cuerpo antes de delegar en \path{lib/typeorm/services/}. Por ultimo, esta arquitectura facilita aplicar controles transversales como el \textit{rate limiting} centralizado en \path{proxy.ts}, la normalizacion de errores mediante \path{lib/api/server.ts} y la reutilizacion de contratos compartidos.

La documentación siguiente se presenta con una estructura similar a una coleccion de Postman: método, URL, autenticación, parámetros o cuerpo de entrada y respuesta principal. Postman organiza las peticiones alrededor de estos elementos de solicitud y respuesta, incluyendo URL, parámetros, autenticación, cabeceras y cuerpo \citep{postmanrequests2026}. En las tablas se omiten cabeceras repetitivas como \texttt{Content-Type: application/json} cuando el cuerpo es JSON; cuando se indica \texttt{multipart/form-data}, la ruta espera subida de archivo.

Como convención general, las rutas autenticadas devuelven \texttt{401} cuando no existe sesión valida y \texttt{403} cuando el rol no está autorizado. Las validaciones de entrada se responden habitualmente con \texttt{400}, los recursos inexistentes con \texttt{404}, los conflictos de datos con \texttt{409} cuando procede y los errores no controlados con \texttt{500}. Todas las rutas bajo \path{/api} pasan por el control de \textit{rate limiting} del \path{proxy.ts}, que puede responder \texttt{429} si se supera la política aplicable.

\newcommand{\apiHeader}{\textbf{Método} & \textbf{Ruta} & \textbf{Acceso} & \textbf{Entrada y respuesta principal}\\}
\newcolumntype{A}[1]{>{\hspace{0pt}\arraybackslash}p{#1}}
\newcommand{\apiTableSetup}{\footnotesize\setlength{\tabcolsep}{3pt}}

\subsubsection{Autenticación, cuenta y perfil}

{\apiTableSetup
\begin{longtable}{A{0.10\textwidth}A{0.27\textwidth}A{0.15\textwidth}A{0.38\textwidth}}
\caption{Endpoints de autenticación, cuenta y perfil}\\
\apiHeader
\endfirsthead
\apiHeader
\endhead
\texttt{GET/POST} & \path{/api/auth/:nextauth*} & Público / sesión & Ruta delegada en \texttt{Auth.js} para CSRF, proveedores, callback de credenciales y consulta de sesión. Respuesta gestionada por la librería.\\
\texttt{POST} & \path{/api/auth/register-request} & Público & JSON con \texttt{email}, \texttt{name}, \texttt{password}, \texttt{company} y \texttt{phone}. Respuesta \texttt{201}: solicitud de alta pendiente de revisión administrativa.\\
\texttt{POST} & \path{/api/auth/logout} & Sesión opcional & Sin cuerpo. Registra el cierre de sesión cuando existe sesión activa y devuelve \texttt{\{ ok: true \}}.\\
\texttt{POST} & \path{/api/account/change-password} & Sesión & \texttt{multipart/form-data} con \texttt{current\_password}, \texttt{new\_password} y \texttt{confirm\_new\_password}. Respuesta: redirección a la pantalla de cambio de contraseña con exito o error.\\
\texttt{PATCH} & \path{/api/profile} & Sesión & JSON con datos de perfil: \texttt{name}, \texttt{email}, \texttt{company}, \texttt{phone}, \path{profile_image_url}, nueva contraseña opcional y \texttt{clientProfile}. Respuesta \texttt{200}: perfil actualizado.\\
\texttt{POST} & \path{/api/profile/upload-image} & Sesión & \texttt{multipart/form-data} con \texttt{file} y \texttt{previousImageUrl}. Respuesta \texttt{200}: \texttt{imageUrl} y \texttt{publicId} de Cloudinary.\\
\end{longtable}
}

\subsubsection{Administración de usuarios y solicitudes}

{\apiTableSetup
\begin{longtable}{A{0.10\textwidth}A{0.27\textwidth}A{0.15\textwidth}A{0.38\textwidth}}
\caption{Endpoints administrativos de usuarios}\\
\apiHeader
\endfirsthead
\apiHeader
\endhead
\texttt{GET} & \path{/api/admin/users} & Admin & Query opcional \texttt{page}, \texttt{pageSize}, \texttt{search}. Respuesta \texttt{200}: listado de usuarios, paginado cuando procede.\\
\texttt{POST} & \path{/api/admin/users} & Admin & JSON \texttt{CreateAdminUserBody}: \texttt{name}, \texttt{email}, \texttt{password}, \texttt{company}, \texttt{phone}, \texttt{roleId}. Respuesta \texttt{201}: usuario creado.\\
\texttt{PATCH} & \path{/api/admin/users/:id} & Admin & JSON de perfil administrativo: datos de usuario, rol y estado. Respuesta \texttt{200}: usuario actualizado.\\
\texttt{PATCH} & \path{/api/admin/users/:id/status} & Admin & JSON con \texttt{statusId}, \texttt{reason} y \texttt{notes}. Respuesta \texttt{200}: estado actualizado y auditado.\\
\texttt{PATCH} & \path{/api/admin/users/:id/role} & Admin & JSON con \texttt{roleId} o \texttt{newRoleId}, \texttt{reason} y \texttt{notes}. Respuesta \texttt{200}: rol actualizado.\\
\texttt{PATCH} & \path{/api/admin/users/:id/password} & Admin & JSON con \texttt{newPassword}, \texttt{reason} y \texttt{notes}. Respuesta \texttt{200}: contraseña cambiada.\\
\texttt{POST} & \path{/api/admin/users/:id/remove} & Admin & Sin cuerpo. Desactiva el usuario indicado y responde con redirección \texttt{303} al listado, o JSON de error si no puede desactivarse.\\
\texttt{POST} & \path{/api/admin/register-user} & Admin & JSON \texttt{RegisterAdminUserBody}: \texttt{name}, \texttt{email}, \texttt{password}, \texttt{type}, datos de contacto y comercial opcional. Respuesta \texttt{201}: usuario registrado.\\
\texttt{GET} & \path{/api/admin/user-requests} & Admin & Sin parámetros. Respuesta \texttt{200}: solicitudes de alta recibidas.\\
\texttt{GET} & \path{/api/admin/user-requests/:id} & Admin & Parámetro de ruta \texttt{id}. Respuesta \texttt{200}: detalle de la solicitud.\\
\texttt{POST} & \path{/api/admin/user-requests/:id/approve} & Admin & JSON opcional con \texttt{commercialId}. Respuesta \texttt{200}: solicitud aprobada y usuario creado o activado.\\
\texttt{POST} & \path{/api/admin/user-requests/:id/reject} & Admin & JSON, formulario o \texttt{multipart/form-data} con \texttt{reason}. Respuesta \texttt{200}: solicitud rechazada.\\
\end{longtable}
}

\subsubsection{Clientes profesionales y asignaciones comerciales}

{\apiTableSetup
\begin{longtable}{A{0.10\textwidth}A{0.27\textwidth}A{0.15\textwidth}A{0.38\textwidth}}
\caption{Endpoints de clientes profesionales y asignación comercial}\\
\apiHeader
\endfirsthead
\apiHeader
\endhead
\texttt{GET} & \path{/api/admin/clients} & Admin & Sin parámetros. Respuesta \texttt{200}: clientes profesionales.\\
\texttt{POST} & \path{/api/admin/clients} & Admin & JSON \texttt{CreateClientRequestBody}: \texttt{name}, \texttt{contactName}, \texttt{taxId}, dirección, ciudad, usuario asociado y notas. Respuesta \texttt{201}: cliente creado.\\
\texttt{GET} & \path{/api/clients/:id} & Sesión con acceso & Parámetro de ruta \texttt{id}. Admin y comercial pueden leer; cliente solo su propia ficha. Respuesta \texttt{200}: ficha del cliente.\\
\texttt{PATCH} & \path{/api/clients/:id} & Admin o cliente propietario & JSON \texttt{UpdateClientRequestBody}: datos de ficha, coordenadas, ventana de visita y notas. Respuesta \texttt{200}: cliente actualizado.\\
\texttt{GET} & \path{/api/clients/:id/geocode} & Admin o cliente propietario & Query \texttt{q} con dirección libre. Respuesta \texttt{200}: coordenadas sugeridas por geocodificación.\\
\texttt{GET} & \path{/api/admin/client-commercial-assignments} & Admin & Query \texttt{clientId} o \texttt{commercialId}. Respuesta \texttt{200}: asignación activa o listado de asignaciones.\\
\texttt{POST} & \path{/api/admin/client-commercial-assignments} & Admin & JSON con \texttt{mode} igual a \texttt{assign}, \texttt{reassign} o \texttt{unassign}, \texttt{clientId}, \texttt{commercialId} y \texttt{notes}. Respuesta \texttt{200} o \texttt{201}: asignación resultante.\\
\texttt{GET} & \path{/api/admin/commercials} & Admin & Sin parámetros. Respuesta \texttt{200}: perfiles comerciales disponibles.\\
\texttt{POST} & \path{/api/admin/commercials} & Admin & JSON \texttt{AdminUpsertCommercialProfileBody}: usuario, código empleado, territorio, jornada, duracion de visitas y puntos de ruta. Respuesta \texttt{200}: perfil comercial creado o actualizado.\\
\texttt{GET} & \path{/api/commercial/clients} & Comercial & Sin parámetros. Respuesta \texttt{200}: clientes asignados al comercial autenticado.\\
\texttt{GET} & \path{/api/commercial/clients/:id} & Comercial & Parámetro \texttt{id}; se valida que el cliente pertenezca al comercial. Respuesta \texttt{200}: detalle del cliente.\\
\texttt{GET} & \path{/api/clients/delivery-estimate} & Cliente & Sin parámetros. Respuesta \texttt{200}: estimación de llegada del reparto planificado para el cliente autenticado.\\
\end{longtable}
}

\subsubsection{Visitas, rutas y perfil comercial}

{\apiTableSetup
\begin{longtable}{A{0.10\textwidth}A{0.27\textwidth}A{0.15\textwidth}A{0.38\textwidth}}
\caption{Endpoints de visitas comerciales, rutas y configuración operativa}\\
\apiHeader
\endfirsthead
\apiHeader
\endhead
\texttt{GET} & \path{/api/admin/commercial-visits} & Admin & Query \texttt{clientId}. Respuesta \texttt{200}: visitas asociadas al cliente.\\
\texttt{POST} & \path{/api/admin/commercial-visits} & Admin & JSON \texttt{CreateCommercialVisitBody}: \texttt{clientId}, \texttt{commercialId}, fecha, tipo, notas y pedidos vinculados. Respuesta \texttt{201}: visita creada y aviso generado al cliente cuando procede.\\
\texttt{GET} & \path{/api/admin/commercial-visits/:id} & Admin & Parámetro \texttt{id}. Respuesta \texttt{200}: detalle de visita y pedidos relacionados.\\
\texttt{PATCH} & \path{/api/admin/commercial-visits/:id} & Admin & JSON \texttt{UpdateCommercialVisitBody}: fecha, tipo, estado, notas, resultado, QR entregados y pedidos. Respuesta \texttt{200}: visita actualizada; si una visita aplazada vuelve a planificada con nueva fecha, se avisa al cliente.\\
\texttt{POST} & \path{/api/admin/commercial-routes} & Admin & JSON con \texttt{commercialId}, \texttt{date} y \texttt{name}. Respuesta \texttt{201}: ruta comercial creada.\\
\texttt{POST} & \path{/api/admin/commercial-routes/add-visit} & Admin & JSON con \texttt{routeId}, \texttt{visitId} y \texttt{order}. Respuesta \texttt{201}: visita agregada a la ruta.\\
\texttt{GET} & \path{/api/commercial/profile} & Comercial & Sin parámetros. Respuesta \texttt{200}: configuración de jornada, duracion y puntos de ruta del comercial.\\
\texttt{PATCH} & \path{/api/commercial/profile} & Comercial & JSON \texttt{UpdateCommercialProfileBody}: jornada, duraciones, direcciones y coordenadas de inicio/fin. Respuesta \texttt{200}: perfil actualizado.\\
\texttt{GET} & \path{/api/commercial/visits} & Comercial & Query opcional \texttt{clientId}, \texttt{dateFrom}, \texttt{dateTo}, \texttt{statusId}, \texttt{visitTypeId}. Respuesta \texttt{200}: visitas filtradas del comercial.\\
\texttt{POST} & \path{/api/commercial/visits} & Comercial & JSON \texttt{CreateCommercialVisitBody}; el comercial queda limitado a sus clientes asignados. Respuesta \texttt{201}: visita creada y aviso generado al cliente cuando procede.\\
\texttt{GET} & \path{/api/commercial/visits/:id} & Comercial & Parámetro \texttt{id}; acceso limitado al comercial propietario. Respuesta \texttt{200}: detalle de visita.\\
\texttt{PATCH} & \path{/api/commercial/visits/:id} & Comercial & JSON \texttt{UpdateCommercialVisitBody}: estado, resultado, QR de entrega y pedidos vinculados. Respuesta \texttt{200}: visita actualizada; si se reubica una visita aplazada, se avisa al cliente.\\
\texttt{GET} & \path{/api/commercial/routes/preview} & Comercial & Query opcional \texttt{startLat}, \texttt{startLng}. Respuesta \texttt{200}: puntos de ruta, ETAs, incidencias y resumen temporal del día.\\
\end{longtable}
}

\subsubsection{Catálogo, productos, recursos y coloración}

{\apiTableSetup
\begin{longtable}{A{0.10\textwidth}A{0.27\textwidth}A{0.15\textwidth}A{0.38\textwidth}}
\caption{Endpoints administrativos de catálogo}\\
\apiHeader
\endfirsthead
\apiHeader
\endhead
\texttt{GET} & \path{/api/admin/catalog/lookups} & Admin & Sin parámetros. Respuesta \texttt{200}: estados de producto y tipos de recurso de apoyo.\\
\texttt{POST} & \path{/api/admin/catalog/upload-image} & Admin & \texttt{multipart/form-data} con \texttt{file} y \texttt{previousImageUrl}. Respuesta \texttt{200}: URL de imagen subida a Cloudinary.\\
\texttt{GET} & \path{/api/admin/catalog/product-categories} & Admin & Query opcional \texttt{search}. Respuesta \texttt{200}: categorías.\\
\texttt{POST} & \path{/api/admin/catalog/product-categories} & Admin & JSON con \texttt{name}, \texttt{description}, \texttt{displayOrder}. Respuesta \texttt{201}: categoría creada.\\
\texttt{GET} & \path{/api/admin/catalog/product-categories/:id} & Admin & Parámetro \texttt{id}. Respuesta \texttt{200}: categoría.\\
\texttt{PATCH} & \path{/api/admin/catalog/product-categories/:id} & Admin & JSON de categoría. Respuesta \texttt{200}: categoría actualizada.\\
\texttt{GET} & \path{/api/admin/catalog/product-lines} & Admin & Query \texttt{productCategoryId}, \texttt{search}. Respuesta \texttt{200}: líneas comerciales.\\
\texttt{POST} & \path{/api/admin/catalog/product-lines} & Admin & JSON con \texttt{name}, \texttt{description}, \texttt{productCategoryId}, \texttt{imageUrl}, \texttt{displayOrder}. Respuesta \texttt{201}: línea creada.\\
\texttt{GET} & \path{/api/admin/catalog/product-lines/:id} & Admin & Parámetro \texttt{id}. Respuesta \texttt{200}: línea comercial.\\
\texttt{PATCH} & \path{/api/admin/catalog/product-lines/:id} & Admin & JSON de línea comercial. Respuesta \texttt{200}: línea actualizada.\\
\texttt{GET} & \path{/api/admin/catalog/product-subcategories} & Admin & Query \texttt{search}, \texttt{productLineId}. Respuesta \texttt{200}: subcategorías.\\
\texttt{POST} & \path{/api/admin/catalog/product-subcategories} & Admin & JSON con \texttt{name}, \texttt{description}, \texttt{productLineId}, \texttt{parentSubcategoryId}, \texttt{imageUrl}, \texttt{displayOrder}. Respuesta \texttt{201}: subcategoría creada.\\
\texttt{GET} & \path{/api/admin/catalog/product-subcategories/:id} & Admin & Parámetro \texttt{id}. Respuesta \texttt{200}: subcategoría.\\
\texttt{PATCH} & \path{/api/admin/catalog/product-subcategories/:id} & Admin & JSON de subcategoría. Respuesta \texttt{200}: subcategoría actualizada.\\
\texttt{GET} & \path{/api/admin/catalog/products} & Admin & Query \texttt{search}, \texttt{productCategoryId}, \texttt{productLineId}, \texttt{productSubcategoryId}, \texttt{statusId}. Respuesta \texttt{200}: productos.\\
\texttt{POST} & \path{/api/admin/catalog/products} & Admin & JSON con nombre, referencia, categoría, línea, subcategoría, imagen, formato, embalaje, información técnica, estado, precio y proveedor. Respuesta \texttt{201}: producto creado.\\
\texttt{GET} & \path{/api/admin/catalog/products/:id} & Admin & Parámetro \texttt{id}. Respuesta \texttt{200}: producto.\\
\texttt{PATCH} & \path{/api/admin/catalog/products/:id} & Admin & JSON de producto. Respuesta \texttt{200}: producto actualizado.\\
\texttt{GET} & \path{/api/admin/catalog/support-resources} & Admin & Query \texttt{search}, \texttt{resourceTypeId}, \texttt{productId}, \texttt{productLineId}. Respuesta \texttt{200}: recursos de apoyo.\\
\texttt{POST} & \path{/api/admin/catalog/support-resources} & Admin & JSON con \texttt{title}, \texttt{description}, \texttt{resourceTypeId}, \texttt{resourceUrl}, \texttt{productId}, \texttt{productLineId}. Respuesta \texttt{201}: recurso creado.\\
\texttt{GET} & \path{/api/admin/catalog/support-resources/:id} & Admin & Parámetro \texttt{id}. Respuesta \texttt{200}: recurso de apoyo.\\
\texttt{PATCH} & \path{/api/admin/catalog/support-resources/:id} & Admin & JSON de recurso de apoyo. Respuesta \texttt{200}: recurso actualizado.\\
\texttt{GET} & \path{/api/admin/catalog/color-charts} & Admin & Query \texttt{productLineId}, \texttt{search}. Respuesta \texttt{200}: cartas de color.\\
\texttt{POST} & \path{/api/admin/catalog/color-charts} & Admin & JSON con \texttt{name}, \texttt{description}, \texttt{productLineId}, \texttt{imageUrl}. Respuesta \texttt{201}: carta creada.\\
\texttt{GET} & \path{/api/admin/catalog/color-charts/:id} & Admin & Parámetro \texttt{id}. Respuesta \texttt{200}: carta de color.\\
\texttt{PATCH} & \path{/api/admin/catalog/color-charts/:id} & Admin & JSON de carta de color. Respuesta \texttt{200}: carta actualizada.\\
\texttt{GET} & \path{/api/admin/catalog/color-references} & Admin & Query \texttt{colorChartId}, \texttt{search}. Respuesta \texttt{200}: referencias cromáticas.\\
\texttt{POST} & \path{/api/admin/catalog/color-references} & Admin & JSON con \texttt{colorChartId}, \texttt{code}, \texttt{name}, descripción, imagen, miniatura y orden. Respuesta \texttt{201}: referencia creada.\\
\texttt{GET} & \path{/api/admin/catalog/color-references/:id} & Admin & Parámetro \texttt{id}. Respuesta \texttt{200}: referencia cromática.\\
\texttt{PATCH} & \path{/api/admin/catalog/color-references/:id} & Admin & JSON de referencia cromática. Respuesta \texttt{200}: referencia actualizada.\\
\end{longtable}
}

\subsubsection{Pedidos, borradores, reparto y cobro}

{\apiTableSetup
\begin{longtable}{A{0.10\textwidth}A{0.27\textwidth}A{0.15\textwidth}A{0.38\textwidth}}
\caption{Endpoints de pedidos y cobros}\\
\apiHeader
\endfirsthead
\apiHeader
\endhead
\texttt{GET} & \path{/api/clients/orders} & Cliente & Sin parámetros. Respuesta \texttt{200}: historial de pedidos del cliente autenticado.\\
\texttt{POST} & \path{/api/clients/orders} & Cliente & JSON con \texttt{notes} y \texttt{lines} de producto, referencia de color y cantidad. Respuesta \texttt{201}: pedido creado.\\
\texttt{GET} & \path{/api/clients/orders/:id} & Cliente & Parámetro \texttt{id}. Respuesta \texttt{200}: detalle del pedido del cliente.\\
\texttt{GET} & \path{/api/clients/orders/draft} & Cliente & Sin parámetros. Respuesta \texttt{200}: borrador activo del cliente.\\
\texttt{PUT} & \path{/api/clients/orders/draft} & Cliente & JSON con \texttt{notes} y \texttt{lines}. Respuesta \texttt{200}: borrador guardado.\\
\texttt{DELETE} & \path{/api/clients/orders/draft} & Cliente & Sin cuerpo. Respuesta \texttt{200}: borrador eliminado.\\
\texttt{POST} & \path{/api/clients/orders/draft/items} & Cliente & JSON con \texttt{productId}, \texttt{colorReferenceId}, \texttt{quantity}. Respuesta \texttt{200}: línea agregada al borrador.\\
\texttt{GET} & \path{/api/commercial/orders} & Comercial & Query opcional \texttt{clientId}, \texttt{pendingDeliveryOnly}. Respuesta \texttt{200}: pedidos visibles para el comercial.\\
\texttt{POST} & \path{/api/commercial/orders} & Comercial & JSON con \texttt{clientId}, \texttt{notes} y \texttt{lines}. Respuesta \texttt{201}: pedido creado para cliente asignado.\\
\texttt{GET} & \path{/api/commercial/orders/:id} & Comercial & Parámetro \texttt{id}. Respuesta \texttt{200}: detalle del pedido accesible para el comercial.\\
\texttt{PATCH} & \path{/api/commercial/orders/:id} & Comercial & JSON con \texttt{statusId}, \texttt{paymentStatusId}, \texttt{paymentMethod}, \texttt{paymentNotes}. Respuesta \texttt{200}: pedido, reparto o cobro actualizado.\\
\texttt{GET} & \path{/api/commercial/orders/draft} & Comercial & Query \texttt{clientId}. Respuesta \texttt{200}: borrador del comercial para ese cliente.\\
\texttt{PUT} & \path{/api/commercial/orders/draft} & Comercial & JSON con \texttt{clientId}, \texttt{notes} y \texttt{lines}. Respuesta \texttt{200}: borrador guardado.\\
\texttt{DELETE} & \path{/api/commercial/orders/draft} & Comercial & Query \texttt{clientId}. Respuesta \texttt{200}: borrador eliminado.\\
\texttt{POST} & \path{/api/commercial/orders/draft/items} & Comercial & JSON con \texttt{clientId}, \texttt{productId}, \texttt{colorReferenceId}, \texttt{quantity}. Respuesta \texttt{200}: línea agregada.\\
\texttt{GET} & \path{/api/admin/orders/:id} & Admin & Parámetro \texttt{id}. Respuesta \texttt{200}: detalle administrativo del pedido.\\
\texttt{PATCH} & \path{/api/admin/orders/:id} & Admin & JSON con \texttt{statusId}, \texttt{paymentStatusId}, \texttt{paymentMethod}, \texttt{paymentNotes}. Respuesta \texttt{200}: estado de pedido o cobro actualizado.\\
\end{longtable}
}

\subsubsection{Historial técnico de salón}

{\apiTableSetup
\begin{longtable}{A{0.10\textwidth}A{0.27\textwidth}A{0.15\textwidth}A{0.38\textwidth}}
\caption{Endpoints del historial técnico de salón}\\
\apiHeader
\endfirsthead
\apiHeader
\endhead
\texttt{GET} & \path{/api/clients/salon-clients} & Cliente & Sin parámetros. Respuesta \texttt{200}: clientes finales del salón.\\
\texttt{POST} & \path{/api/clients/salon-clients} & Cliente & JSON con \texttt{name}, \texttt{phone}, \texttt{email}, \texttt{notes}. Respuesta \texttt{201}: cliente final creado.\\
\texttt{GET} & \path{/api/clients/salon-clients/:id} & Cliente & Parámetro \texttt{id}. Respuesta \texttt{200}: ficha, servicios y sugerencias.\\
\texttt{PATCH} & \path{/api/clients/salon-clients/:id} & Cliente & JSON de cliente final. Respuesta \texttt{200}: ficha actualizada.\\
\texttt{POST} & \path{/api/clients/salon-clients/:id/services} & Cliente & JSON con fecha, tipo de servicio, notas, resultado, descripción técnica, formula, productos usados e imágenes. Respuesta \texttt{201}: servicio creado.\\
\texttt{PATCH} & \path{/api/clients/salon-clients/:id/services/:serviceId} & Cliente & JSON de servicio técnico. Respuesta \texttt{200}: servicio actualizado.\\
\texttt{DELETE} & \path{/api/clients/salon-clients/:id/services/:serviceId} & Cliente & Sin cuerpo. Respuesta \texttt{200}: servicio eliminado.\\
\texttt{GET} & \path{/api/clients/salon-clients/:id/services/:serviceId/technical-email} & Cliente & Parámetros \texttt{id} y \texttt{serviceId}. Respuesta \texttt{200}: borrador de correo técnico editable.\\
\texttt{GET} & \path{/api/clients/salon-service-templates} & Cliente & Sin parámetros. Respuesta \texttt{200}: plantillas de servicio del cliente.\\
\texttt{POST} & \path{/api/clients/salon-service-templates} & Cliente & JSON con nombre, tipo, notas, resultado, descripción técnica, formula, notas técnicas y productos. Respuesta \texttt{201}: plantilla creada.\\
\texttt{DELETE} & \path{/api/clients/salon-service-templates/:templateId} & Cliente & Parámetro \texttt{templateId}. Respuesta \texttt{200}: plantilla eliminada.\\
\texttt{POST} & \path{/api/clients/salon-service-result-images} & Cliente & \texttt{multipart/form-data} con \texttt{file}. Respuesta \texttt{200}: imagen de resultado subida a Cloudinary.\\
\texttt{DELETE} & \path{/api/clients/salon-service-result-images} & Cliente & JSON con \texttt{imageUrl}. Respuesta \texttt{200}: imagen de resultado eliminada.\\
\end{longtable}
}

\subsubsection{Comunicaciones internas}

{\apiTableSetup
\begin{longtable}{A{0.10\textwidth}A{0.27\textwidth}A{0.15\textwidth}A{0.38\textwidth}}
\caption{Endpoints de comunicaciones, promociones, formaciones y recordatorios}\\
\apiHeader
\endfirsthead
\apiHeader
\endhead
\texttt{GET} & \path{/api/admin/communications/segments} & Admin & Query opcional \texttt{search}. Respuesta \texttt{200}: segmentos de cliente.\\
\texttt{POST} & \path{/api/admin/communications/segments} & Admin & JSON con \texttt{code}, \texttt{name}, \texttt{description}, \texttt{criteria}. Respuesta \texttt{201}: segmento creado.\\
\texttt{GET} & \path{/api/admin/communications/segments/:id} & Admin & Parámetro \texttt{id}. Respuesta \texttt{200}: segmento.\\
\texttt{PATCH} & \path{/api/admin/communications/segments/:id} & Admin & JSON de segmento. Respuesta \texttt{200}: segmento actualizado.\\
\texttt{DELETE} & \path{/api/admin/communications/segments/:id} & Admin & Sin cuerpo. Respuesta \texttt{200}: segmento eliminado.\\
\texttt{GET} & \path{/api/admin/communications/client-segments} & Admin & Query opcional \texttt{search}. Respuesta \texttt{200}: asignaciones cliente-segmento.\\
\texttt{POST} & \path{/api/admin/communications/client-segments} & Admin & JSON con \texttt{clientId}, \texttt{segmentId}, \texttt{notes}. Respuesta \texttt{201}: asignación creada.\\
\texttt{DELETE} & \path{/api/admin/communications/client-segments/:id} & Admin & Parámetro \texttt{id}. Respuesta \texttt{200}: asignación eliminada.\\
\texttt{GET} & \path{/api/admin/communications/promotions} & Admin & Query opcional \texttt{search}. Respuesta \texttt{200}: promociones administrativas.\\
\texttt{POST} & \path{/api/admin/communications/promotions} & Admin & JSON con titulo, descripción, tipo, beneficio, fechas, estado, alcance por producto, línea, cliente o segmento y \texttt{deliveryChannels}. Respuesta \texttt{201}: promoción creada y notificaciones generadas por los canales seleccionados cuando queda activa.\\
\texttt{GET} & \path{/api/admin/communications/promotions/:id} & Admin & Parámetro \texttt{id}. Respuesta \texttt{200}: promoción.\\
\texttt{PATCH} & \path{/api/admin/communications/promotions/:id} & Admin & JSON de promoción, incluyendo \texttt{deliveryChannels} opcional. Respuesta \texttt{200}: promoción actualizada y notificaciones emitidas si pasa a estar activa.\\
\texttt{DELETE} & \path{/api/admin/communications/promotions/:id} & Admin & Sin cuerpo. Respuesta \texttt{200}: promoción eliminada.\\
\texttt{GET} & \path{/api/admin/communications/trainings} & Admin & Query opcional \texttt{search}. Respuesta \texttt{200}: formaciones administrativas.\\
\texttt{POST} & \path{/api/admin/communications/trainings} & Admin & JSON con titulo, descripción, fecha, ubicacion, modalidad, contenido, estado, capacidad y \texttt{deliveryChannels}. Respuesta \texttt{201}: formación creada y notificaciones generadas por los canales seleccionados cuando queda publicada.\\
\texttt{GET} & \path{/api/admin/communications/trainings/:id} & Admin & Parámetro \texttt{id}. Respuesta \texttt{200}: formación.\\
\texttt{PATCH} & \path{/api/admin/communications/trainings/:id} & Admin & JSON de formación, incluyendo \texttt{deliveryChannels} opcional. Respuesta \texttt{200}: formación actualizada y notificaciones emitidas si pasa a estar publicada.\\
\texttt{DELETE} & \path{/api/admin/communications/trainings/:id} & Admin & Sin cuerpo. Respuesta \texttt{200}: formación eliminada.\\
\texttt{GET} & \path{/api/communications/promotions} & Cliente y comercial & Sin parámetros. Respuesta \texttt{200}: promociones visibles para el usuario autenticado.\\
\texttt{GET} & \path{/api/communications/trainings} & Cliente y comercial & Sin parámetros. Respuesta \texttt{200}: formaciones publicadas para el usuario.\\
\texttt{POST} & \path{/api/communications/trainings/:id/enrollment} & Cliente y comercial & JSON opcional con \texttt{notes}. Respuesta \texttt{201}: inscripción registrada.\\
\texttt{DELETE} & \path{/api/communications/trainings/:id/enrollment} & Cliente y comercial & Sin cuerpo. Respuesta \texttt{200}: inscripción cancelada.\\
\texttt{GET} & \path{/api/communications/notifications} & Cliente y comercial & Sin parámetros. Respuesta \texttt{200}: notificaciones del usuario; para clientes sincroniza antes los recordatorios internos de visitas planificadas para el día actual.\\
\texttt{PATCH} & \path{/api/communications/notifications} & Cliente y comercial & Sin cuerpo. Respuesta \texttt{200}: todas las notificaciones marcadas como leídas.\\
\texttt{PATCH} & \path{/api/communications/notifications/:id} & Cliente y comercial & Parámetro \texttt{id}. Respuesta \texttt{200}: notificación marcada como leída.\\
\texttt{GET} & \path{/api/communications/push-subscriptions/public-key} & Cliente y comercial & Sin parámetros. Respuesta \texttt{200}: \texttt{configured} y clave pública VAPID disponible para suscribir el navegador.\\
\texttt{POST} & \path{/api/communications/push-subscriptions} & Cliente y comercial & JSON \texttt{PushSubscription}: \texttt{endpoint}, \texttt{expirationTime}, \texttt{keys.p256dh} y \texttt{keys.auth}. Respuesta \texttt{201}: suscripcion creada o actualizada y contador de suscripciones activas.\\
\texttt{DELETE} & \path{/api/communications/push-subscriptions} & Cliente y comercial & JSON opcional con \texttt{endpoint}. Respuesta \texttt{200}: suscripcion revocada y contador de suscripciones activas.\\
\texttt{GET} & \path{/api/communications/reminders} & Cliente y comercial & Sin parámetros. Respuesta \texttt{200}: recordatorios propios.\\
\texttt{POST} & \path{/api/communications/reminders} & Cliente y comercial & JSON con \texttt{title}, \texttt{body}, \texttt{scheduledAt}, \texttt{status}. Respuesta \texttt{201}: recordatorio creado.\\
\texttt{PATCH} & \path{/api/communications/reminders/:id} & Cliente y comercial & JSON con \texttt{status}. Respuesta \texttt{200}: estado del recordatorio actualizado.\\
\end{longtable}
}

\subsubsection{Configuración transversal}

{\apiTableSetup
\begin{longtable}{A{0.10\textwidth}A{0.27\textwidth}A{0.15\textwidth}A{0.38\textwidth}}
\caption{Endpoints transversales de configuración}\\
\apiHeader
\endfirsthead
\apiHeader
\endhead
\texttt{GET} & \path{/api/admin/operations/status} & Admin & Sin parámetros. Respuesta \texttt{200}: resumen técnico agregado de operación transversal de \texttt{M7}, usado por scripts y validación interna.\\
\texttt{GET} & \path{/api/admin/operations/report} & Admin & Sin parámetros. Respuesta \texttt{200}: fichero \texttt{CSV} descargable con el informe técnico operativo de \texttt{M7}.\\
\texttt{GET} & \path{/api/admin/enterprise-operations} & Admin & Sin parámetros. Respuesta \texttt{200}: configuraciones, integraciones externas, operaciones recientes y propuestas a proveedor de \texttt{M7}, como base para integración futura.\\
\texttt{POST} & \path{/api/admin/enterprise-operations/operations} & Admin & JSON con \texttt{integrationId}, \texttt{operationType}, \texttt{dataType}, \texttt{status} y \texttt{result}. Respuesta \texttt{201}: operación de integración registrada.\\
\texttt{POST} & \path{/api/admin/enterprise-operations/supplier-proposals} & Admin & JSON con \texttt{productIds}, \texttt{quantity}, \texttt{reason}, \texttt{notes} y \texttt{status}. Respuesta \texttt{201}: propuesta de pedido a proveedor generada con sus líneas.\\
\texttt{GET} & \path{/api/admin/audit/export} & Admin & Query \texttt{type} igual a \texttt{access} o \texttt{management}, \texttt{days} y \texttt{limit} opcionales. Respuesta \texttt{200}: fichero \texttt{CSV} descargable con registros recientes de auditoría.\\
\texttt{GET} & \path{/api/admin/integrations/status} & Admin & Sin parámetros. Respuesta \texttt{200}: resumen e inventario técnico de integraciones de soporte.\\
\texttt{GET} & \path{/api/admin/support/status} & Admin & Sin parámetros. Respuesta \texttt{200}: resumen e inventario técnico de soporte, PWA y compatibilidad.\\
\texttt{GET} & \path{/api/admin/settings/rate-limits/status} & Admin técnico & Sin parámetros. Respuesta \texttt{200}: diagnóstico agregado de políticas de rate limiting cuando la ruta técnica está habilitada por entorno; en caso contrario devuelve \texttt{404}.\\
\texttt{GET} & \path{/api/admin/settings/rate-limits} & Admin técnico & Sin parámetros. Respuesta \texttt{200}: políticas de \textit{rate limiting}, valores por defecto y valores configurados cuando la ruta técnica está habilitada; en caso contrario devuelve \texttt{404}.\\
\texttt{PUT} & \path{/api/admin/settings/rate-limits} & Admin técnico & JSON con \texttt{policies[]}: \texttt{name}, \texttt{enabled}, \texttt{maxRequests}, \texttt{windowMinutes}. Respuesta \texttt{200}: políticas actualizadas solo con la ruta técnica habilitada.\\
\end{longtable}
}
